\documentclass[11pt]{article}
\usepackage[a4paper,margin=1in]{geometry}
\usepackage{amsmath,amssymb}
\usepackage{booktabs,longtable}
\usepackage{graphicx}
\usepackage[hyphens]{url}
\usepackage[colorlinks=true,linkcolor=black,urlcolor=blue,citecolor=blue]{hyperref}
\setlength{\parskip}{0.55em}\setlength{\parindent}{0pt}
\providecommand{\tightlist}{\setlength{\itemsep}{0.2em}\setlength{\parskip}{0pt}}

\title{Cover Letter — Paper E1: resubmission to Fisheries Research}
\author{\textbf{Amin Abaee}\\ Independent Researcher\\[2pt] ORCID: \href{https://orcid.org/0000-0002-0019-1842}{0000-0002-0019-1842}\\ Email: \href{mailto:amin_abaee@ut.ac.ir}{amin_abaee@ut.ac.ir}}
\date{September 18, 2026}
\begin{document}
\maketitle

\section{\texorpdfstring{Cover Letter --- Paper E1 (resubmission to
\emph{Fisheries
Research})}{Cover Letter --- Paper E1 (resubmission to Fisheries Research)}}\label{cover-letter-paper-e1-resubmission-to-fisheries-research}

\textbf{Title:} Forecasting biomass under structural non-stationarity:
an out-of-sample evaluation of surplus-production models for Northern
cod (Gadus morhua) \textbf{Author:} Amin Abaee (Independent Researcher;
ORCID 0000-0002-0019-1842; amin\_abaee@ut.ac.ir) \textbf{Journal:}
\emph{Fisheries Research} \textbf{Date:} September 18, 2026

\begin{center}\rule{0.5\linewidth}{0.5pt}\end{center}

Dear Editor,

I resubmit the manuscript \textbf{``Forecasting biomass under structural
non-stationarity: an out-of-sample evaluation of surplus-production
models for Northern cod''} for consideration in \emph{Fisheries
Research}, following your initial decision. The revision responds
directly to the three concerns raised at that stage; I have addressed
two by clarifying and specifying, and the third --- the decisive one ---
by running the demanded test empirically.

\subsection{What the paper intends to
be}\label{what-the-paper-intends-to-be}

The paper is a \textbf{forecast-comparison study}, not a proposed
assessment method and not a sustainability verdict. Its intention is
narrowly stated: whether adding structure to a surplus-production ladder
improves out-of-sample forecasts of the quantity that management advice
actually consumes --- spawning stock biomass --- is an empirical
question with a stated burden of proof. The manuscript tests that
question on the canonical demanding case (Northern cod, NAFO 2J3KL)
under a scoring rule fixed before the scores were read, with persistence
treated as an active competitor, not a scaling denominator. The claims
are deliberately conservative: the negative result is scoped to the
fitted estimator family and ladder, and the paper explicitly does not
conclude anything about the stock's sustainability or about the adequacy
of historical management decisions.

\subsection{Response to the three
concerns}\label{response-to-the-three-concerns}

\textbf{Concern 1 --- Technical and computational specification.} We
understood the first concern as: the ladder could not be independently
reconstructed from the manuscript. The revised paper carries a dedicated
\textbf{Appendix A (A.1--A.7)} that consolidates the full specification
in one place: the map equations, per-window estimation mechanics
(bounds, multi-start, the frozen autoregressive rule, catch handling),
the rolling-origin machinery and metrics, the executable retention rule,
and the uncertainty layer (Diebold--Mariano with Newey--West HAC;
moving-block bootstrap; 20,000 replications, seed 0). Together with the
reproducibility package in Supplementary SI-6 and the deterministic
campaign scripts archived with the revision, every number in the results
is regenerable byte-for-byte.

\textbf{Concern 2 --- Scope and status of the benchmark.} We understood
the second concern as a misreading risk: that the manuscript might be
taken to propose a new assessment method for the stock. The opening
section now states the scope in its first paragraph: the ladder is a
scored test of structural additions against a demanding naive baseline,
in the tradition of persistent-benchmark diagnostics (MASE; Hyndman and
Koehler, 2006) that assessment science already treats as canonical.
Nothing in the paper is offered as a replacement assessment tool, and no
management advice is issued from it.

\textbf{Concern 3 --- The choice of scoreboard.} We understood the third
--- and in our view correct --- concern as: the verdict might depend on
the smoothed assessment reconstruction as target rather than on the
models themselves; a raw monitoring series could tell a different story.
The first revision could only reserve this question for follow-up.
\textbf{The revision now answers it empirically.} Section 3.8 scores the
same frozen ladder, under the same pre-registered rule, against the raw
autumn research-vessel survey index (1983--2015, 33 consecutive years;
Schijns et al., 2021, Table 3), on origin sets identical to the
assessment runs. The retained set is empty on the monitoring target as
well: persistence records 120.5 kt-equivalents at one year against
156.3--371.3 for the structural ladder, and 249.2 at five years against
355.5--880.5 --- the persistence advantage \emph{increases} on the raw
index (+29.8\% vs +17.1\% at one year; +42.7\% vs +9.0\% at five years).
Section 4.3 explains the mechanism explicitly: smoothing a fixed target
and substituting a different, noisier target are different
interventions; observation noise compounds through iterated structural
forecasts while persistence pays only the bounded one-step noise floor.
The conversion constant, the one non-separating cell (the stock-flow
module at one year, whose gap remains a deficit either way), and the
verdict-level invariance of every outcome to that constant, are all
documented in the text and SI-7. The headline therefore changed form:
\textbf{no surplus-production module is retained on either a
reconstructed or a monitoring target.}

\subsection{What changed since the first
submission}\label{what-changed-since-the-first-submission}

\begin{enumerate}
\def\labelenumi{\arabic{enumi}.}
\tightlist
\item
  \textbf{Scope clarifications} in §1 and the abstract (concern 2), with
  the paper's intention stated before any method.
\item
  \textbf{Appendix A.1--A.7} technical/computational specification
  (concern 1); SI extended to v5, including the reproducibility
  inventory.
\item
  \textbf{A new empirical object}: the pre-registered survey-index
  campaign (§3.8, Table 12b; design frozen before any survey score was
  read), the target-contrast mechanism discussion (§4.3), and the
  machinery-fidelity, determinism, and scale-sensitivity verification
  recorded in SI-7.
\end{enumerate}

The manuscript is self-contained; all inputs, scripts, and result files
are archived (https://github.com/MIKEAA2020/general-sustainability). The
work is original, is not under consideration elsewhere, and I have no
conflicts of interest.

I believe the revision has converted an arguably \emph{target-dependent}
negative result into a \emph{target-robust} one, and that this is now
squarely within the journal's remit. Thank you for your consideration.

Yours sincerely,

\textbf{Amin Abaee} Independent Researcher ORCID: 0000-0002-0019-1842
Email: amin\_abaee@ut.ac.ir

\end{document}
